% =====================================================================
% Section 4 — Implementation
% Facts verified against the codebase (JxMVC.Core, v3.4.0):
%   pipeline: MainLxServlet; rate limiter: JxRateLimiter.estimatedUsage;
%   OIDC: JxOAuth.verifyJwtClaims; tests: JxTestSuite (custom, no JUnit).
% =====================================================================
\section{Implementation}\label{sec:impl}

This section highlights the implementation techniques that are most
relevant to the framework's claims of being small, secure-by-default,
and dependency-free. All code paths described here are in the core
module (\texttt{JxMVC.Core}, v3.4.0), which comprises 54 classes and
compiles against a single \texttt{provided}-scope dependency, the
Jakarta~EE~10 API.

\subsection{Front controller and routing}

The entire pipeline of Section~\ref{sec:pipeline} is implemented in a
single front-controller servlet (\texttt{MainLxServlet}) whose
\texttt{service} method applies the security headers and conditional
gzip that must wrap \emph{every} response, and whose internal dispatch
runs the ordered stages. Routing (\texttt{BaseDispatcher}) resolves a
request through three mechanisms in decreasing specificity: explicit
annotated routes (\texttt{@JxGetMapping} and siblings), path templates
with \texttt{\{var\}} capture, and a
\texttt{/controller/action/arg} convention. The router returns the
correct \texttt{405 Method Not Allowed} when a path matches but the
verb does not, and treats \texttt{HEAD} as \texttt{GET}. To keep
reflection off the hot path, annotated-method metadata is resolved once
per controller class and reused across requests.

\subsection{Rate limiting: a sliding-window counter}

Rate limiting (\texttt{JxRateLimiter}) is keyed by
\texttt{VERB:controller\#action} rather than by URL, so an attacker
cannot dilute a limit by rotating a path variable. It uses a
sliding-window \emph{counter} (two time-aligned buckets, current and
previous) rather than a per-request timestamp log, trading exactness
for $O(1)$ memory per client. The estimated request count over the
window is
\[
  \hat{n} \;=\; c_{\text{prev}} \cdot \Bigl(1 - \tfrac{t}{W}\Bigr) \;+\; c_{\text{cur}},
\]
where $W$ is the window length, $t$ the elapsed time into the current
window, and $c_{\text{prev}}, c_{\text{cur}}$ the previous and current
bucket counts. This is the standard sliding-window approximation, here
implemented with the JDK alone.

\subsection{OIDC without libraries}

The OAuth~2.0/OpenID~Connect client (\texttt{JxOAuth}) implements the
authorization-code flow with PKCE using only \texttt{java.security},
\texttt{java.math}, and \texttt{java.util.Base64}. It generates an
anti-CSRF \texttt{state} and a PKCE verifier, sends the S256 challenge
$\text{code\_challenge} = \text{BASE64URL}(\text{SHA-256}(\text{verifier}))$,
and exchanges the code for tokens. Crucially, it does not merely decode
the \texttt{id\_token}: it \textbf{verifies the RS256 signature}
against the provider's JWKS. The JWKS is fetched and cached (one-hour
TTL); the RSA public key is reconstructed from its modulus and exponent
via \texttt{KeyFactory}/\texttt{RSAPublicKeySpec}; and the signature is
checked with \texttt{Signature.getInstance("SHA256withRSA")}. The
verifier then validates the \texttt{iss}, \texttt{aud}, and \texttt{exp}
claims (with a 60-second clock skew) and \textbf{accepts only
\texttt{RS256}}---any other algorithm, including \texttt{none}, is
rejected with \texttt{401}. Implementing token verification, rather
than trusting decoded claims, closes a well-known class of OIDC
vulnerabilities, and doing so with the JDK alone demonstrates that a
zero-dependency posture is compatible with correct, modern
authentication.

\subsection{Defence-in-depth details}

Several protections are worth noting as concrete implementations rather
than configuration switches. Password hashing (\texttt{JxPasswords})
uses \texttt{PBKDF2WithHmacSHA256} with a 16-byte random salt and a
self-describing encoding (\texttt{pbkdf2\$sha256\$iters\$salt\$hash}),
verified in constant time. File uploads reject a fixed blocklist of
dangerous extensions (e.g.\ \texttt{.jsp}, \texttt{.class},
\texttt{.war}, \texttt{.sh}, \texttt{.exe}, \texttt{.php}) \emph{even
when} a wildlist is configured, and normalise paths to prevent
traversal. The JSON parser (\texttt{JxJson}) caps input at
\SI{10}{\mega\byte} and nesting depth at 64, returning \texttt{400}
rather than exhausting memory on adversarial input. Redirects to
external hosts are blocked unless explicitly enabled.

\subsection{A dependency-free test suite}

Consistent with the zero-dependency principle, the framework is tested
\emph{without} JUnit or any external testing library. A custom runner
(\texttt{JxTestSuite}) executes fifteen test classes via reflection,
aggregating pass/fail counts and exiting non-zero on any failure, and a
lightweight helper constructs mock request/response objects so
controllers can be tested without a servlet container. The suite
comprises \NTESTS{} checks and deliberately targets the critical
production paths---connection pooling under concurrency, gzip
passthrough, cache stampede, concurrent WebSocket broadcast, JSON error
escaping, path traversal, and mass assignment---rather than trivial
getters, so that ``dependency-free'' extends even to how the framework
proves its own correctness.
